\section{各级标题与段落格式测试}
本章主要用于测试一级、二级、三级标题的格式规范，以及普通正文的段落排版。

根据规范要求，正文部分的各种中文基本都应该是宋体、小四号，而数字和英文应该是 Times New Roman、小四号体。整篇文章采用首行缩进两个汉字，采用 1.25 倍的行距，并且段前段后距均为 0 磅。

所有的单数页页眉为当前章节标题，而偶数页页眉固定为“上海电力大学硕士/博士学位论文”。我们可以通过翻到下一页来验证偶数页页眉。

\subsection{二级标题的规范检查}
这是二级标题。按照要求，二级标题应该左对齐顶格，小三号宋体并且加黑，1.25倍行距，段前空一行，段后空 0.5 行。

我们在此处多写几段文字，以便将内容推送到第二页验证偶数页页眉是否正常显示。

这是另起的一段。所有的行文必须满足左右两端对齐的要求。这一页是奇数页，如果排版准确，页眉处的文字应该是“第一章 各级标题与段落格式测试”。

\subsubsection{三级标题测试}
这部分是三级标题的测试区。按照要求，三级标题应当左起空两个汉字字符，使用四号宋体并且加黑。段前应空 0.5 行，段后空 0 行。请仔细检查本标题左侧是否有相当于两个汉字宽度的留白。
